\section{Where the active ingredients may lie}\label{sec:ingredients}
A defender of CBT has an obvious reply: patients do not restructure during a panic; they practice when calm, so that the response becomes automatic. We think this is the strongest objection to the executive-function argument, and we accept its force. It also changes what is being claimed. If repeated rehearsal during calm periods does the work, the active ingredient is practice and habit formation, not in-the-moment reasoning, and the relevant question is how much executive effort each rehearsal costs.

A second observation points the same way. Exposure, which is central to CBT for anxiety, is behavioural: the patient approaches the feared stimulus and learns, from experience, that the feared outcome does not occur. Its logic resembles extinction learning more than debate. If most of the benefit of the CBT package comes from exposure and skills practice, the argument against \emph{restructuring} as the key ingredient is strengthened without any claim that CBT does not work.
\citegap{dismantling trials (Jacobson et al.), component meta-analyses, inhibitory-learning accounts of exposure, homework compliance meta-analyses}
